\section{The problem and the shorter route}
\label{rh-purpose}

The Riemann hypothesis asserts that every zero of the Riemann zeta
function in the open strip $0<\Re s<1$ has real part $1/2$. This
manuscript studies a shorter route to that assertion, within a broader
reconstruction of Levent Alp\"oge's announced proof of the Riemann
hypothesis for Dirichlet $L$-functions \cite{rh-alpoge,rh-companion}.
The arithmetic construction that would complete the argument has not
been supplied here. The purpose is to make its necessary output smaller
and its mathematical dependencies explicit.

There are three principal steps. First, an upper bound for one real
prime sum excludes zeros to the right of the critical line. Second,
positivity and finite character averaging bound that prime sum by a
single arithmetic discrepancy, three auxiliary quadratic sums, and
explicit errors. Third, geometric degree estimates and the trace errors of
integral maps provide sufficient conditions for that discrepancy.
The first step is complete. The second is a complete implication from
the stated arithmetic identity and sampled estimates. The third contains
proved comparison theorems, with substantial construction inputs still
to be established on the same family.

The arithmetic comparisons have a common short structure. Integral
linearized equations lead to a determinant valuation; a finite slice
with at most two points bounds its multiplicity. On the geometric side,
two filtrations control the finite determinant order, and one elliptic
curve with complex multiplication supplies its exact period value.
The original model system and ordinary families are now constructed.
The remaining work concerns their precise section and support
identifications, uniform integral bounds and global analytic estimates.

The distinction between a comparison theorem and its arithmetic input is
essential. For example, proving that an isotropic family of matrices
has good rank properties does not construct integral maps with those
matrices, metrics and heights. Similarly, a bound for a logarithmic
derivative does not establish a lower bound for an arithmetic
intersection number. We retain those identifications as hypotheses
where they enter.

\subsection{Constructed inputs and the precise arithmetic frontier}
\label{shared-current-arithmetic-frontier}

The analytic endpoint is an implication from a compatible arithmetic
estimate. Several of its original construction inputs are now supplied
in the detailed construction proofs\cite{current-construction-record}.
We distinguish those inputs from the quantitative assertions still
needed to apply the endpoint. All constructions below retain the
original fields, integral lattices, maps, measures and local conditions.

\paragraph{Models and sections.}
The integral PEL models, coordinated level changes, formal arcs and
global dominating model give the original complete model system.
The normalized complex theta dictionary and local integral charts give
the prescribed good cover. The actual raw theta section exists over
every coefficient algebra, has its literal classical value and lies in
the required bounded cusp space in the stated positive-weight range.
One clearing integer, chosen before the generator, supplies its
integral image with the proved denominator bound. This does not yet
give short integral spanning families. The actual automorphic lines
are now nef and big in the required gap range; effective adjoint
separation supplies the prescribed birational cusp systems. The
section-count argument still needs short integral spanning families,
the prescribed height comparisons and its numerical remainder.

\paragraph{Selmer and automorphic constructions.}
The ordinary automorphic modules, integral Hecke determinant and
compact types are constructed at the original field and level.
The Taylor--Wiles comparisons and ordinary families retain the exact
integral equations, multiplier and torsion. Propagation from a supplied
stable support prime is proved. The actual seed now transfers to the
original normalized deformation space with the same integral Hecke
quotient and the required starting-chart dimension. Access to the
required characteristic-prime components, two-sided global rigidity
and the uniform weighted Selmer product remain distinct obligations. Characteristic-zero points alone do
not control the characteristic-prime fibre.

\paragraph{Theta identities and global norms.}
The literal Gaussian is integrable on the full prescribed slab and
reproduces the cusp space with its positive scalar, including the
even endpoint $n=2r$. For the actual Hosts, the exact Hermitian
Siegel--Weil identity holds in its stronger range $n>4r$, with the
original level markings, automorphism masses and Weil phases.
The harmonic heat comparison transfers that zero-Jacobi identity
to the invariant derivatives used by the kernel argument. The finite
Schwartz system, level-theta span and exact identity-cover transfer
now give the actual algebraic raw-theta span under the original
level and weight hypotheses. Short integral spanning bounds remain
quantitative requirements. Arithmetic folding, finite phase matching
and the complete framed-return sum are also supplied at each fixed
setting under its original gap. Identification with the actual period
Gram operators and sufficiently small trace errors still require
arithmetic work; finite setting-dependent constants are not uniform bounds.

\paragraph{The downstream acceptance condition.}
For each accuracy, the constructions must coexist at every entry of
the chosen fourth-power tables. Their signed arithmetic cost, including
rank, integral-index, height and projection losses, must satisfy the
stated subpower bound on arbitrarily late intervals. In the RH route,
four surviving characters simplify the projection but impose no bound
on the auxiliary field degree or geometric dimension. The construction
frontier is therefore substantive in both routes. Its complete
quantifiers are recorded with the final arithmetic budget; none is
replaced by a fixed-rank or characteristic-zero assertion.


\subsection{A corrected prime sum}

All logarithms are natural. Let $\Lambda(p^j)=\log p$ for a prime $p$
and integer $j\ge1$, and let $\Lambda(n)=0$ otherwise. Write $\gamma$
for Euler's constant. For $x\ge1$ set
\begin{equation}\label{rh-G}
 G(x)=\sum_{n\le x}\Lambda(n)(x/n-1)-x\log x+(1+\gamma)x.
\end{equation}
The summand vanishes when $x=n$, so $G$ is continuous. Its correction
removes the pole of $\zeta$ at one from the Mellin transform. The
one-sided target is
\begin{equation}\label{rh-simple-target}
 G(x)\le C_\epsilon x^{1/2+\epsilon}
       \qquad(x\ge X_\epsilon),
\end{equation}
for every $\epsilon>0$, with finite constants depending on $\epsilon$.
No corresponding lower bound is required.

Auxiliary characters enter only in the arithmetic comparison. For a
primitive nonprincipal Dirichlet character $\chi$, define
\begin{equation}\label{rh-Gchi}
 G_\chi(x)=\sum_{n\le x}\Lambda(n)\chi(n)(x/n-1)
                   +x\frac{L'(1,\chi)}{L(1,\chi)},
 \qquad G_1=G.
\end{equation}
The symbol $1$ means the character of conductor one. Principal
characters written at larger moduli have different values at the
removed primes and must not be substituted into this definition.
Products of characters below are always replaced by their primitive
inducing characters when $G_\chi$ is formed.

We use the classical analytic continuation and functional equations
of $\zeta$ and Dirichlet $L$-functions, the Euler product in $\Re s>1$,
and $L(1,\chi)\ne0$ for nonprincipal $\chi$; see
\cite{rh-dlmf-zeta,rh-dlmf-dirichlet}. In particular,
$\zeta(s)=(s-1)^{-1}+\gamma+O(s-1)$ near one. These are background
analytic inputs. The newer uniform zero-free estimates, when invoked,
will be stated separately and will not be identified with RH.

\subsection{The two useful acceptance conditions}

The eventual condition \eqref{rh-simple-target} has a short proof and
is the easiest endpoint to understand. For construction purposes, a
weaker condition is useful. Put $z_+=\max(z,0)$ and, for fixed $h>0$,
\[
 \mathcal A_hF(x)=\frac1h\int_0^hF(xe^v)\,dv,
 \qquad \mathcal A_0F=F.
\]
One may instead bound
\[
 \int_Y^{Y^\beta}(\mathcal A_hG(x))_+
                          x^{-\theta}\,\frac{dx}{x}
\]
by a quantity smaller than every positive power of $Y$, on arbitrarily
large intervals, for exponents $\theta$ tending to $1/2$.
At each accuracy, $h\ge0$, $\beta>1$ and $\theta$ are fixed while the
intervals move to infinity. This permits signed cancellation between
nearby scales and between terms in the final arithmetic budget. The
positive part is taken only after those signed terms are combined.

Sections~\ref{rh-analytic} and \ref{rh-transfer} prove the analytic
endpoint and the direct transfer. Sections~\ref{rh-geometry} and
\ref{rh-energy} develop the smaller geometric and matrix conditions.
Section~\ref{rh-budget} states the complete remaining target and shows
how the quantifiers and parameter losses fit together.
